
\begin{proof}[\pfref{prop:greedy}]
  We prove the result by induction. Fix $x\in\cX$, and let
  $\ystar_1,\ldots,\ystar_H\ldef{}\ystar(x)$. Fix $h\in\brk{H}$, and
  assume by induction that $\yhat_{h'}=\ystar_{h'}$ for all $h'<h$. We
  claim that in this case,
  \begin{align}
    \pi_h(\ystar_h\mid{}\yhat_{1},\ldots,\yhat_{h-1},x)
    = \pi_h(\ystar_h\mid{}\ystar_{1},\ldots,\ystar_{h-1},x)>1/2, 
  \end{align}
  which implies that $\yhat_h=\ystar_h$. To see this, we observe that
  by Bayes' rule,
  \begin{align}
    \label{eq:4}
    \pi(\ystar_1,\ldots,\ystar_H\mid{}x)
    &\leq{}     \pi(\ystar_1,\ldots,\ystar_h\mid{}x)\\
    &=\prod_{h'=1}^{h}\pi_{h'}(\ystar_{h'}\mid{}\ystar_1,\ldots,\ystar_{h'-1},x)
    \leq{}\pi_h(\ystar_h\mid{}\ystar_{1},\ldots,\ystar_{h-1},x).
  \end{align}
  If we were to have
  $\pi_h(\ystar_h\mid{}\yhat_{1},\ldots,\yhat_{h-1},x)=\pi_h(\ystar_h\mid{}\ystar_{1},\ldots,\ystar_{h-1},x)\leq{}1/2$,
  it would contradict the assumption that
  $\pi(\ystar_1,\ldots,\ystar_H\mid{}x)>1/2$. This proves the result.
  \end{proof}

